% ============================================
% Chapter: Conclusions, Limitations, and Future Work
% ============================================

\chapter{Conclusions, Limitations, and Future Work}
\label{ch:conclusions}

Conclusions,
Limitations, and
Future Work



\section{Introduction to the Chapter}


This chapter closes the GreenLab project report by evaluating the system against the
eight objectives defined in Section 1.4, consolidating the limitations documented
throughout Chapters 3 through 5 into a single account, and proposing a concrete
direction for continued development. Chapter 5 traced individual functional and nonfunctional requirements to their implementation evidence (Section 5.6) and recorded
seven specific implementation deviations as they were discovered (Section 5.8). This
chapter operates at the level of the project as a whole: it evaluates whether GreenLab
achieved what it set out to achieve, states its limitations plainly, and proposes the
roadmap that would carry it from a demonstrated prototype to a validated deployment.


\section{Achievement of Project Objectives}


Table 6.1 evaluates each objective from Section 1.4 against the evidence accumulated
across Chapters 3--5, principally the requirements traceability matrix in Section 5.6.

Table 6.1: Objective-by-objective achievement status.

# Objective Status Evidence
The hardware, AI, backend, and dashboard
tiers communicate exclusively through the
defined REST interfaces described in Section 4.1
Layered four-tier
and exercised in Chapter 5: the ESP32 firmware
architecture with
Achieved never bypasses the backend to reach the AI
clear interface
service or dashboard directly, and each tier's
contracts
internal implementation (§5.2--5.5) can be
replaced without altering another tier's
contract.

Sections 3.5 and 3.6 specify 135 requirements
Functional and non- across nine functional areas and six non-
functional functional quality dimensions. The Section 5.6
Achieved
requirements in traceability matrix confirms every sampled
formal IEEE 830 style requirement traces to implementation evidence
or an explicitly recorded gap.


# Objective Status Evidence
Multimodal
The Section 5.7 worked trace demonstrates HC-
occupancy detection
SR04, DHT11, and LDR readings combining with
combining
Achieved the YOLO service's tracked person count within
embedded sensors
a single automation decision (FR-OD-01, FR-
and YOLO-based
OD-02: Implemented).
vision

The LDR-zone × occupancy-classification matrix
Lighting control via a (Table 4.6) is implemented as a data-driven rule
Achieved
decision matrix lookup, satisfying the auditability and
configurability goal stated in Objective 4.

The temperature-zone × occupancy-
Ventilation control
classification table (Table 4.7) follows the same
via a fan activation Achieved
pattern as the lighting matrix and is exercised
table
directly in the Section 5.7 worked example.

FR-AC-01 through FR-AC-04 are marked
Secure ASP.NET Core Implemented in the Section 5.6 traceability
REST API with matrix, supported by JWT-based authentication
authentication, Achieved (§4.7.1), a three-role access model (§4.7.2), and
RBAC, and audit an append-only audit log (§4.7.3), each
logging demonstrated against concrete
request/response evidence in Chapter 5.

Sections 5.4--5.5 document eight distinct
dashboard views spanning all four required
capabilities, with FR-MD-01 and FR-MD-03 both
Angular dashboard
marked Implemented. One presentational
for monitoring,
Achieved inconsistency --- a third device-state pill
override, scheduling,
rendered for a currently-unwired relay channel
and administration
--- is recorded as a known limitation in Section
\section{and does not affect the four core capabilities}

this objective requires.

Complete, print- Achieved The report's figures, tables, and formatting
ready project report by this satisfy Objective 8. The report was previously

# Objective Status Evidence
incomplete against its own six-chapter
structure (no Chapter 6) and carried an
revision unresolved citation-numbering placeholder in
Chapter 2; both are resolved as part of this
revision.


Seven of the eight objectives were fully achieved by the delivered system with no
qualification beyond the implementation-level deviations recorded in Section 5.8.
Objective 8, concerning the report itself, is brought to completion by this chapter and the
Chapter 2 correction above.


\section{Summary of Project Contribution}


GreenLab set out to close a specific gap identified in the literature review (Section 2.11):
no reviewed source combined occupancy-aware automation with role-based governance,
authenticated override, and audit accountability in a single system designed for the
operational realities of a shared academic laboratory. The Table 2.1 comparison matrix
positions GreenLab as the only reviewed system offering multimodal sensing, decisionmatrix-driven automation, and a governed, auditable control layer simultaneously. The
as-built system demonstrates that this combination is achievable within a four-tier
architecture where each tier owns a narrow, well-defined responsibility, connected
through authenticated REST interfaces rather than direct point-to-point coupling. The
Section 5.7 worked example demonstrates the complete path from a physical entry event
to an audited, authorised lighting decision, executing in under 600 milliseconds end to
end, supporting the practical feasibility of the design.


\section{Limitations}


\subsection{Implementation Deviations}


Section 5.8 recorded seven deviations between the Chapter 3/4 specification and the asbuilt system as they were discovered during implementation. They are restated here as
formal project limitations.

Table 6.2: Implementation deviations carried forward as project limitations.




Limitation Requirement Nature of deviation
affected
Firmware control
Changing a lighting or temperature threshold
thresholds are
NFR-MAIN-01 requires reflashing the device rather than
compile-time
calling a runtime configuration endpoint.
constants

The backend's response contract models four
Two of four modelled relay channels (lightOn, fanOn, relay3On,
Architecture
relay channels are relay4On); only the first two are physically
(§4.2, §5.8)
unwired connected to GPIO output in the current
prototype.

On connectivity loss, the firmware reLocal network-
evaluates the LDR reading against a local
fallback logic re-
FR-FH-02 threshold rather than holding the last
evaluates rather than
backend-commanded relay state, as
freezes
originally specified.

This role was introduced during
The Teaching implementation to reflect a real
FR-AC-02, FRAssistant role is organisational need but does not appear in
AC-03
unmodelled the Chapter 3 use case model or stakeholder
analysis.

The ``occupancy
This alert type is implemented and functions
outside session'' alert Fault handling
correctly but was not anticipated by the
has no formal (§5.5.2)
Chapter 3 functional requirement set.
requirement ID

labId and deviceId are integer literals in the
Device identity is
firmware rather than provisioned at
compiled into NFR-MAIN-01
runtime, meaning each new device
firmware
deployment requires a firmware rebuild.

Inference sampling
NFR-PERF-03 FRAME_SKIP=3 yields a variable inference
cadence is frame-
rate depending on the camera's actual frame
count-based, not


Limitation Requirement Nature of deviation
affected
time-based rate rather than a fixed elapsed-time interval.


\subsection{Project-Level Limitations}


Beyond the seven implementation-level deviations above, the project carries limitations
that apply to its overall scope and maturity rather than to any single component:

 No field deployment or empirical validation. As stated in the project scope
(Section 1.6), GreenLab has not been installed in a live laboratory environment,
and no energy savings, detection accuracy, or latency figures are reported as
measured GreenLab results beyond the single traced example in Section 5.7.

 No multi-user or extended-duration testing. The system has been exercised
through individual workflow traces and integration checkpoints (Section 1.5), but
not through sustained multi-day operation or concurrent multi-instructor usage.

 Single-prototype hardware scale. All implementation evidence in Chapter 5
reflects one ESP32 unit and one camera source; behaviour under multiple
concurrent laboratory deployments against a shared backend has not been
exercised.


\section{Future Work}


The following four-phase roadmap is proposed as the natural continuation of the project.

Phase 1 --- Configuration and provisioning (near-term). Replace compile-time
firmware thresholds and compiled-in device identity with a runtime configuration and
provisioning workflow: a /api/devices/{id}/config endpoint fetched at boot, and a QR-code
or BLE commissioning step for labId/deviceId assignment. This resolves NFR-MAIN-01
without requiring firmware rebuilds for calibration or deployment changes.

Phase 2 --- Requirements and behavioural closure (near-term). Formally add the
Teaching Assistant role to the Chapter 3 use case model and stakeholder analysis, specify
the ``occupancy outside session'' alert as FR-FH-04, replace the frame-count-based
FRAME_SKIP sampling gate with a time-based one (resolving NFR-PERF-03), and correct
the network-fallback logic to hold rather than re-evaluate the last commanded relay state
(resolving FR-FH-02). These four items close every specification-level gap recorded in
Section 5.8 without requiring new architecture.

Phase 3 --- Hardware completion (mid-term). Wire the two currently-unused relay
channels to additional laboratory circuits, bringing the physical hardware in line with
the four-channel contract already modelled in the backend response schema, and
expand the hardware node to a small pilot batch (5--10 units) to validate the provisioning
workflow from Phase 1 under realistic multi-device conditions.

Phase 4 --- Field pilot and empirical evaluation (long-term). Deploy GreenLab in a
defined set of pilot laboratories and collect the measurements explicitly excluded from
this report's scope: detection accuracy under real lighting and occupancy conditions,
measured energy consumption before and after deployment, and instructor-reported
usability feedback. This phase converts the project's design-time claims into empirically
supported results and is the single highest-value step toward a production-ready system.


\section{Concluding Remarks}


GreenLab demonstrates that occupancy-aware laboratory automation and operational
governance are complementary layers of the same system rather than competing
concerns: the same event that triggers an automatic lighting decision is, in the same code
path, written to an immutable audit record and surfaced on an administrator's
dashboard. Seven of the project's eight objectives were fully achieved by the delivered
system, with the eighth --- completion of the report itself --- closed by this chapter. The
limitations recorded in Section 6.4 mark the specific, named boundary between what has
been built and demonstrated on a development network, and what remains to be proven
through field deployment. The four-phase roadmap in Section 6.5 defines the path from
that boundary to a validated, production-ready deployment.


\section{Chapter Summary}


This chapter evaluated GreenLab against its eight stated objectives (Section 6.2), finding
seven fully achieved on technical evidence traced through Chapters 3--5 and the eighth
completed by this revision of the report. Section 6.3 restated the project's core
contribution: the combination of multimodal occupancy sensing with a governed,
auditable control layer, a combination absent from the literature reviewed in Chapter 2.
Section 6.4 consolidated seven implementation-level deviations and three project-level
limitations into a single honest account of what has, and has not, been demonstrated.
Section 6.5 proposed a four-phase roadmap --- configuration and provisioning,
requirements and behavioural closure, hardware completion, and field pilot evaluation
--- as the concrete path from the current prototype to a validated deployment. Together



with the corrected Chapter 2 transition, this chapter completes the six-chapter structure
of the GreenLab graduation project report.




